%% Copyright (c) 2018-2020 Weitian LI <wt@liwt.net>
%% CC BY 4.0 License
%%
%% Created: 2018-04-11
%%

% Chinese version
\documentclass[zh]{resume}

% File information shown at the footer of the last page
\fileinfo{
  % \faCopyright{} 2018--2020, Weitian LI \hspace{0.5em}
  % \creativecommons{by}{4.0} \hspace{0.5em}
  % \githublink{liweitianux}{resume} \hspace{0.5em}
  \faEdit{} \today
}

\name{德生}{杨}

% \keywords{BSD, Linux, Programming, Python, C, Shell, DevOps, SysAdmin}

% \tagline{\texorpdfstring{\icon{\faBinoculars} }{}<position-to-look-for>}
% \tagline{<current-position>}

% \photo[
%   shape=<circular|square>,    % default is circular
%   position=<left|right>,      % default is left
% ]{<width>}{<filename>}
% Example:
% \photo[shape=square]{5.5em}{photo}

\profile{
  \mobile{187-7975-0679}
  \email{desyang@qq.com}
  \github{desyang-hub} \\
  \university{南昌航空大学}
  \degree{计算机科学与技术 \textbullet 硕士}
  \birthday{2000-09}
  \address{南昌}
  % Custom information:
  % \icontext{<icon>}{<text>}
  % \iconlink{<icon>}{<link>}{<text>}
}

\begin{document}
\makeheader

%======================================================================
% Summary & Objectives
%======================================================================
\begin{abstract}
物理学专业（射电天文方向）直博研究生，有扎实的物理、数学与统计学基础，
擅长数据建模与分析，热衷计算机和网络技术，
有 10 年的 Linux 和 BSD 使用经验，熟练掌握 Shell、Python 和 C 语言编程。
积极实践自由开源精神，
在 \link{https://github.com/liweitianux}{GitHub} 上分享多个项目，
是 \link{https://www.dragonflybsd.org}{DragonFly BSD} 操作系统的开发者，
并积极参与其他多个开源项目。
\end{abstract}

%======================================================================
\sectionTitle{求职意向}{\faBullseye}
%======================================================================
\begin{itemize}
  \item C/C++开发工程师
\end{itemize}

%======================================================================
\sectionTitle{技能和语言}{\faWrench}
%======================================================================
\begin{competences}
  \comptence{操作系统}{%
    \icon{\faLinux} Linux （10 年），
    \icon{\faFreebsd} DragonFly BSD \& FreeBSD （7 年）
  }
  \comptence{编程}{%
    Python, C, Shell, R, Tcl/Tk
  }
  \comptence{工具}{%
    SSH, Git, Make, Tmux, Vi, Ansible
  }
  \comptence{数据分析}{%
    R, Pandas; Matplotlib, ggplot2; Keras, Scikit-learn
  }
  \comptence{网站开发}{%
    Flask, JavaScript, jQuery, Bootstrap
  }
  \comptence{\icon{\faLanguage} 语言}{
    \textbf{英语} --- 读写（优良），听说（日常交流）
  }
\end{competences}

%======================================================================
\sectionTitle{教育背景}{\faGraduationCap}
%======================================================================
\begin{educations}
  \education%
    {至今}%
    [2024.09]%
    {南昌航空大学}%
    {信息工程学院}%
    {计算机科学与技术}%
    {硕士在读}

  \separator{0.5ex}
  \education%
    {2024.06}%
    [2020.09]%
    {南昌航空大学}%
    {信息工程学院}%
    {计算机科学与技术}%
    {学士}
\end{educations}

%======================================================================
\sectionTitle{计算机技能}{\faCogs}
%======================================================================
\begin{itemize}
  \item DragonFly BSD 操作系统开发者：
    200+ 代码提交；内核以及系统工具；
    在邮件列表和 IRC 频道交流和回答问题
  \item 使用 Ansible 管理 VPS，部署个人域名邮箱、权威 DNS、网站、Git、IRC 等服务
  \item 搭建并管理课题组的工作站、计算集群（4 节点）和网络设备
  \item 参与配置和测试上海天文台的 SKA 高性能计算集群原型机
    （1 管理节点 + 1 存储节点 + 4 计算节点）
  \item 设计并开发了\enquote{2014 第一届中国—新西兰联合 SKA 暑期学校}的整个网站
    （Django, Bootstrap, jQuery）
\end{itemize}

%======================================================================
\sectionTitle{个人项目}{\faCode}
%======================================================================
\begin{itemize}
  \item \link{https://github.com/liweitianux/atoolbox}{\texttt{atoolbox}}:
    (Python, Shell)
    多年来累积的各种工具，帮助管理系统、执行常用任务、分析天文数据等
  \item \link{https://github.com/liweitianux/dfly-update}{\texttt{dfly-update}}:
    (Shell)
    DragonFly BSD 系统更新程序
  \item \link{https://github.com/liweitianux/openrcs}{\texttt{openrcs}}:
    (C)
    改进 \texttt{OpenBSD RCS}，使其与 \texttt{GNU RCS} 足够兼容
  \item \link{https://github.com/liweitianux/fg21sim}{\texttt{fg21sim}}:
    (Python)
    模拟低频射电天空图像
  \item \link{https://github.com/liweitianux/cdae-eor}{\texttt{cdae-eor}}:
    (Python, Keras)
    使用卷积去噪自动编码器（CDAE）分离宇宙再电离（EoR）信号
  \item \link{https://github.com/liweitianux/chandra-acis-analysis}{\texttt{chandra-acis-analysis}}:
    (Python, Shell, Tcl)
    X~射线天文观测数据的半自动化分析程序
  \item \link{https://github.com/liweitianux/resume}{\texttt{resume}}:
    (\LaTeX)
    \emph{此简历}的模板和源文件
\end{itemize}

%======================================================================
\sectionTitle{科研成果}{\faAtom}
%======================================================================

\begin{itemize}
  \setlength{\itemsep}{0.4ex}
  \renewcommand{\labelitemi}{\color{gray}\tiny$\bullet$}
  
  % \item \textbf{《基于深度学习的图像分割方法》} \\
  %       \textit{IEEE Transactions on Image Processing} (第一作者) \hfill 2025 
  \item \textbf{发明专利：一种矩形印章文字识别方法} (第一作者) \hfill 2024(已授权) \\
      专利号：ZL 2024 1 0373918.4 
  \item \textbf{发明专利：一种文档图像表格类型判定方法} (第二作者,导师一作) \hfill 2025(已授权) \\
      专利号：ZL 2025 1 0865310.8 
  \item \textbf{发明专利：一种轻量级表格结构识别方法} (第三作者) \hfill 2025(已授权)  \\
      专利号：ZL 2025 1 0409433.0     
\end{itemize}

%======================================================================
\sectionTitle{实习经历}{\faBriefcase}
%======================================================================
\begin{experiences}
  \experience%
    [2018.04]%
    {2018.08}%
    {数据工程师 @ 上海领脉网络科技（初创公司）}%
    [\begin{itemize}
      \item 从 Amazon 网页搜索并挖取商品与广告信息
        （Python, Requests, BeautifulSoup）
      \item 配置 Airflow 服务器和数据库等基础设施，
        定期从 Amazon 获取产品销售与广告投放等数据
      \item 开发网站（Flask, jQuery），帮助客户优化 Amazon 广告投放
    \end{itemize}]

  \separator{0.5ex}
  \experience%
    [2013.07]%
    {2013.09}%
    {网站开发 @ 97 随访（初创公司）}%
    [\begin{itemize}
      \item 后端开发（Django），完成用户注册、数据存储和搜索等功能
      \item 前端开发（jQuery, AJAX），对患者各项指标随时间的变化进行可视化
    \end{itemize}]
\end{experiences}


%======================================================================
\sectionTitle{荣誉奖项}{\faTrophy}
%======================================================================
\begin{itemize}
  \item 国家励志奖学金 \hfill 2021
  \item 睿抗机器人开发者大赛(RAICOM) CAIP编程设计赛 国赛二等奖 \hfill 2023
  \item 优秀研究生 \hfill 2025
\end{itemize}

\end{document}
